% © 2026 Ing. David Jaroš — CC BY-NC-ND 4.0
%
% This work is licensed under a Creative Commons Attribution-NonCommercial-NoDerivatives
% 4.0 International License (CC BY-NC-ND 4.0).
%
% License History: Earlier drafts (up to v0.3) were released under CC BY 4.0.
% From v0.4 onward, all material is released under CC BY-NC-ND 4.0 to protect
% the integrity of the theoretical work during ongoing academic development.
%
% See LICENSE.md for full license text.

\ifdefined\INCLUDEMODE
\else
  \documentclass[12pt,a4paper]{article}
  \usepackage[a4paper, margin=2.5cm]{geometry}
  \usepackage{amsmath, amssymb, amsthm}
  \usepackage{mathtools}
  \usepackage{hyperref}
  \usepackage{xcolor}
  \usepackage{mdframed}
  \usepackage{booktabs}

  \newtheorem{theorem}{Theorem}[section]
  \newtheorem{lemma}[theorem]{Lemma}
  \newtheorem{proposition}[theorem]{Proposition}
  \newtheorem{conjecture}[theorem]{Conjecture}
  \theoremstyle{definition}
  \newtheorem{definition}[theorem]{Definition}
  \theoremstyle{remark}
  \newtheorem{remark}[theorem]{Remark}

  \newmdenv[backgroundcolor=yellow!20, linecolor=orange, linewidth=1.5pt]{gapbox}
  \newmdenv[backgroundcolor=green!15, linecolor=green!60!black, linewidth=1.5pt]{resultbox}
  \newmdenv[backgroundcolor=red!10, linecolor=red!60, linewidth=1.5pt]{deadendbox}
  \newmdenv[backgroundcolor=blue!8, linecolor=blue!50, linewidth=1.5pt]{insightbox}

  \title{\textbf{Gap A --- Approach A2: Modular Curve Geometry}\\
         \large Is $B_{\mathrm{phenom}}$ the Volume of a Hyperbolic Modular Curve?}
  \author{David Jaro\v{s}}
  \date{2026-03-06}

  \begin{document}
  \maketitle

  \begin{abstract}
  We test whether the phenomenological coefficient $B_{\mathrm{phenom}} \approx 46.3$
  required for $\alpha^{-1} = 137$ in the Unified Biquaternion Theory can be
  identified with the hyperbolic volume of a modular curve $X_0(N)$ at the
  UBT-selected levels $N = 137$ and $N = 139$.  We compute volumes, genera, and
  natural normalisations.  The volume $\mathrm{vol}(X_0(139)) \approx 146.6$ is
  a factor $\approx 3$ too large; dividing by $\pi$ gives $\approx 46.7$, which is
  tantalizingly close.  We assess whether this near-coincidence is accidental or
  structurally significant.  Companion script:
  \texttt{tools/compute\_modular\_curve\_genus.py}.
  \end{abstract}
\fi

%──────────────────────────────────────────────────────────────────────────────
\section{Context: Volume Formula for $X_0(N)$}
\label{sec:context}
%──────────────────────────────────────────────────────────────────────────────

The modular curve $X_0(N)$ is the compactified quotient
$\Gamma_0(N)\backslash\mathfrak{H}^*$ of the upper half-plane.  Its hyperbolic
area (in the normalisation where $\mathfrak{H}/\mathrm{SL}(2,\mathbb{Z})$ has
area $\pi/3$) is
\begin{equation}
  \mathrm{vol}(X_0(N))
  \;=\; \frac{\pi}{3}\,\psi(N),
  \quad
  \psi(N) \;=\; N \prod_{p \mid N} \!\left(1 + \frac{1}{p}\right),
  \label{eq:vol_formula}
\end{equation}
where the product runs over prime divisors of $N$.

\begin{insightbox}
\textbf{Key numbers:}
\begin{align*}
  N = 137 \text{ (prime)}: &\quad
    \psi(137) = 137\cdot\tfrac{138}{137} = 138,
    \quad
    \mathrm{vol}(X_0(137)) = \tfrac{138\pi}{3} = 46\pi \approx 144.5 . \\
  N = 139 \text{ (prime)}: &\quad
    \psi(139) = 139\cdot\tfrac{140}{139} = 140,
    \quad
    \mathrm{vol}(X_0(139)) = \tfrac{140\pi}{3} \approx 146.6 .
\end{align*}
\end{insightbox}

Neither $144.5$ nor $146.6$ equals $B_{\mathrm{phenom}} \approx 46.3$ directly.
The ratio $146.6 / 46.3 \approx 3.17 \approx \pi$, which motivates the test in
Section~\ref{sec:volume_test}.

%──────────────────────────────────────────────────────────────────────────────
\section{Genus of $X_0(N)$ for the UBT Attractor Primes}
\label{sec:genus}
%──────────────────────────────────────────────────────────────────────────────

The genus of $X_0(p)$ for a prime $p$ is given by the Riemann-Hurwitz formula:
\begin{equation}
  g(X_0(p))
  \;=\; \begin{cases}
    \displaystyle \frac{p - 13}{12}
      & \text{if } p \equiv 1 \pmod{12}, \\[6pt]
    \displaystyle \left\lfloor \frac{p+1}{12} \right\rfloor
      & \text{otherwise (with corrections for elliptic points)},
  \end{cases}
  \label{eq:genus_formula}
\end{equation}
with the full formula involving the number of elliptic points of orders 2 and 3
and the number of cusps.  For $p = 137$ and $p = 139$:

\begin{center}
\begin{tabular}{lccccc}
  \toprule
  Prime $p$ & $\mu_0(p) = p+1$ & Elliptic pts & Cusps & $g(X_0(p))$ \\
  \midrule
  $137$ & $138$ & 0 (since $137 \nmid \text{special values}$) & 2 & $11$ \\
  $139$ & $140$ & 0 & 2 & $11$ \\
  \bottomrule
\end{tabular}
\end{center}

Both primes yield the same genus $g = 11$.  Consequently the ratio
$g(X_0(139))/g(X_0(137)) = 1$, which provides no explanation for the ratio
$B_{\mathrm{phenom}}/B_0 \approx 1.842$.

\begin{deadendbox}
\textbf{Dead end (genus ratio):}
$g(X_0(137)) = g(X_0(139)) = 11$.  The genus ratio is exactly~1, not~1.842.
The genus does not explain the gap.
\end{deadendbox}

%──────────────────────────────────────────────────────────────────────────────
\section{Volume Normalisation Tests}
\label{sec:volume_test}
%──────────────────────────────────────────────────────────────────────────────

\subsection{Direct volume comparison}

\begin{center}
\begin{tabular}{lll}
  \toprule
  Quantity & Numerical value & Matches $B_{\mathrm{phenom}} \approx 46.3$? \\
  \midrule
  $\mathrm{vol}(X_0(137)) = 46\pi$ & $144.51$ & No ($\times 3.12$) \\
  $\mathrm{vol}(X_0(139)) = 140\pi/3$ & $146.61$ & No ($\times 3.17$) \\
  $\mathrm{vol}(X_0(137))/\pi = 46$ & $46.00$ & \textbf{Very close!} \\
  $\mathrm{vol}(X_0(139))/\pi = 140/3$ & $46.67$ & Close ($+0.8\%$) \\
  \bottomrule
\end{tabular}
\end{center}

\begin{insightbox}
\textbf{Near-coincidence:}
\[
  \frac{\mathrm{vol}(X_0(137))}{\pi} \;=\; 46 ,
  \qquad
  B_{\mathrm{phenom}} \;\approx\; 46.3.
\]
The difference is $0.3/46 \approx 0.65\%$.  Whether this is an accident or a
structural equality $B_{\mathrm{phenom}} = \mathrm{vol}(X_0(137))/\pi$ is the
central question of this approach.
\end{insightbox}

\subsection{Physical interpretation of the $1/\pi$ normalisation}

If $B_{\mathrm{phenom}} = \mathrm{vol}(X_0(N))/\pi$ for $N = 137$, then one
needs a physical mechanism that produces a $1/\pi$ factor.  Natural candidates
include:
\begin{itemize}
  \item One-loop integration over the modular curve with the standard
        hyperbolic measure divided by the area of the fundamental domain.
  \item The ratio of the UBT path-integral measure to the Poincar\'{e} measure.
  \item A Weil height factor $\pi$ arising from the Weierstrass $\wp$-function
        at the CM point $\tau = i\sqrt{137}$.
\end{itemize}
None of these is yet derived from UBT first principles.

\subsection{Correction: $B_{\mathrm{phenom}}$ vs $\mathrm{vol}(X_0(137))/\pi = 46$}

The observed fine-structure constant gives $\alpha^{-1} = 137.036$, so
$B_{\mathrm{phenom}}$ must reproduce this with sufficient precision:
\begin{equation}
  B_{\mathrm{phenom}}
  \;=\; \frac{2A}{\ln(n^*) + 1}
  \;\approx\; \frac{2A}{\ln 137 + 1}
  \;\approx\; \frac{2A}{5.920}.
\end{equation}
The value $B = 46$ (exactly) would give $n^* = e^{2A/46 - 1}$.  To get $n^* = 137$
one needs $B \approx 46.3$ (within $0.65\%$), which is consistent with
$B = \mathrm{vol}(X_0(137))/\pi = 46$ only if higher-order corrections shift $B$
by $+0.3$.

%──────────────────────────────────────────────────────────────────────────────
\section{Index and Cusp Structure}
\label{sec:index_cusps}
%──────────────────────────────────────────────────────────────────────────────

The index of $\Gamma_0(N)$ in $\mathrm{SL}(2,\mathbb{Z})$ is
$\mu(\Gamma_0(N)) = N\prod_{p\mid N}(1 + 1/p) = \psi(N)$.  For primes:
\begin{equation}
  \mu(\Gamma_0(137)) \;=\; 138 \;=\; 2 \times 3 \times 23
  \;=\; 3 \times 46,
\end{equation}
\begin{equation}
  \mu(\Gamma_0(139)) \;=\; 140 \;=\; 4 \times 5 \times 7.
\end{equation}

The factorisation $138 = 3 \times 46$ is striking: if the one-loop computation
produces a factor of $3$ relative to the volume (from the three spatial
dimensions, for instance), then $\mu(\Gamma_0(137))/3 = 46$ exactly.

\begin{gapbox}
\textbf{Candidate formula:}
\[
  B_{\mathrm{phenom}} \;\stackrel{?}{=}\;
  \frac{\mu(\Gamma_0(137))}{3} \;=\; 46 .
\]
This is $0.65\%$ below the required $46.3$.  The $0.3$ discrepancy could arise
from:
\begin{itemize}
  \item Higher-loop corrections to $B_0 = 8\pi$.
  \item The precise value $n^* = 137.036$ rather than $n^* = 137$ exactly.
  \item An additive correction from the cusp contribution to the volume.
\end{itemize}
\end{gapbox}

%──────────────────────────────────────────────────────────────────────────────
\section{Verdict and Status}
\label{sec:verdict}
%──────────────────────────────────────────────────────────────────────────────

\begin{center}
\begin{tabular}{lll}
  \toprule
  Formula & Value & Status \\
  \midrule
  $\mathrm{vol}(X_0(139)) = 140\pi/3$ & $146.6$ & No direct match \\
  $\mathrm{vol}(X_0(137))/\pi = 46$ & $46.0$ & Near-coincidence ($< 1\%$) \\
  $\mu(\Gamma_0(137))/3 = 46$ & $46.0$ & Same near-coincidence \\
  $g(X_0(137)) = 11$ & $11$ & No direct connection to ratio \\
  \bottomrule
\end{tabular}
\end{center}

\begin{gapbox}
\textbf{Status: OPEN.}  No direct match found between $B_{\mathrm{phenom}}
\approx 46.3$ and modular curve volumes or genera.  However, the near-equality
\[
  B_{\mathrm{phenom}} \;\approx\; \frac{\mathrm{vol}(X_0(137))}{\pi}
  \;=\; 46
\]
(within $0.65\%$) is too close to be immediately dismissed.  Two sub-cases:
\begin{enumerate}
  \item The equality is exact ($B_{\mathrm{phenom}} = 46$ exactly), and the
        fine-structure constant is $\alpha^{-1} \approx 136.97$ with higher
        corrections accounting for the residual $0.036$.
  \item The equality is approximate, and the true value of $B_{\mathrm{phenom}}$
        arises from a different mechanism (Approach A1, A3, or A4) that
        numerically happens to be close to $46$.
\end{enumerate}
Distinguishing these requires higher-loop UBT computations.
See companion script \texttt{tools/compute\_modular\_curve\_genus.py}.
\end{gapbox}

\ifdefined\INCLUDEMODE
\else
  \end{document}
\fi
